\scanpage{20}
à isom près est \uline{indépendant} du choix de la \unsure{fact.|fact\textsuperscript{n}} Notation $R^q_! f_*(F)$.

\[
R^0_! f_*(F)=f_!(F)=\text{image directe à support propre sur }Y.
\]

\deleted{Propriétés} Propriétés plus surprenantes à certains égards que $R^q f_*(F)$, car \uline{aussi bon que si $f$ propre}

(i) \uline{Comparaison avec cas classique} sur $\mathbb C$

(ii) \uline{Th. de changement de base} \deleted{\unsure{cf|cf.|c}}, \uline{dim} \uline{coh.}

(iii) \uline{\deleted{Suite spectrale de Leray}}\inserted{Th. de finitude}

(iv) Suite spectrale de Leray

% slanted
\margin{phénomènes spéciaux}
(v) Si $f$ quasi-fini\inserted{\unsure{sép.}}, alors les $R^i_! f_*(F)=0$

\deleted{D'autres \unsure{faits|points|part} =}
(vi) Si $f$ \uline{étale}, on a une trace can.

\[
\mathrm{Tr}_f(F):f_!(F)\unsure{\mathrel{-}|\longrightarrow}F
\]

[d'où loi covariante, pour $f:X'\mathrel{-}X$ morphisme étale -- p. ex. imm ouverte -- ] de préschémas compactifiables sur $Y$

\[
R^q_!g'_*(F')\mathrel{-}R^q_!g_*(f_!(F'))
\]

ou encore

\[
R^q_!g'_*(\deleted{\unsure{F}}f^*(F))\longrightarrow R^q_!g_*(F)\qquad\big]
\]

(vi) \uline{Suite exacte de cohomologie relative}

% archive sheet number: 19
